\section{Related Work}
\label{sec:related}
\subsection{Machine unlearning in LLMs}
Several methods address unlearning in LLMs. Yao \textit{et al.}~\cite{yao2024large} introduced gradient ascent (GA) on forget samples paired with gradient descent (GD) on retain samples. However, the diversity of LLM corpora makes retain set collection intractable, causing GA to erode utility and risk catastrophic forgetting. To mitigate this, Zhang \textit{et al.}~\cite{zhang2024negative} proposed Negative Preference Optimization (NPO), a Direct Preference Optimization (DPO)-inspired objective that slows GA divergence via adaptive gradient weighting. However, NPO still inherits GA at its core, leaving utility vulnerable to unsampled knowledge erosion. 

Shifting from gradient-based objectives, Li \textit{et al.}~\cite{li2024wmdp} introduced the WMDP benchmark alongside Representation Misdirection for Unlearning (RMU), which steers forget activations toward a random unit vector while anchoring retain activations. However, the resulting models often produce incoherent outputs rather than clean refusals. To eliminate retain data dependence, Wang \textit{et al.}~\cite{wang2024llm} proposed Forget-data-only Loss Adjustment (FLAT), which maximizes f-divergence between template and forget responses using only forget data. However, FLAT operates at the batch level and remains ineffective for single data point removal. Revisiting GA’s update mechanics, Wang \textit{et al.}~\cite{wang2025rethinking} proposed the G-effect diagnostic alongside Weighted Gradient Ascent (WGA), which assigns per-instance importance weights to curb over-unlearning. However, G-effect only measures impacts on observed retain samples, leaving collateral damage on unseen regions undetected. From a different perspective, Zade \textit{et al.}~\cite{zade2026attention} proposed Attention Smoothing Unlearning (ASU), which casts unlearning as self-distillation from a forget-teacher with elevated attention temperature to flatten memorized token associations. However, the dual forward pass doubles GPU memory usage, making it impractical at scale. 

Notably, none of these methods are training-free or designed for single-sample forgetting—both of which are essential for interactive machine unlearning at test time. Our work addresses these two gaps.

\subsection{Test-time training (TTT)}
Since our framework performs unlearning at inference time, we review existing test-time training approaches. Sun \textit{et al.}~\cite{sun2024learning} replace the RNN hidden state with a small model updated via self-supervised gradient descent at each token, achieving linear complexity with transformer-like scaling. However, it only compresses patterns within the current sequence rather than acquiring new knowledge. Extending this idea, Akyurek \textit{et al.}~\cite{akyurek2024surprising} fine-tune models via LoRA at test time using synthetic tasks generated through leave-one-out augmentation. However, synthetic data only approximates the true distribution, and pseudo-label quality degrades on novel tasks. 

At a larger scale, Behrouz \textit{et al.}~\cite{behrouz2024titans} introduced surprise-driven selective memorization with sliding-window attention to scale beyond 2M context. However, Titans still memorize contextual patterns rather than acquiring genuinely new knowledge from interactions. Targeting attention, Bansal \textit{et al.}~\cite{bansal2025let} proposed qTTT, which applies gradient updates to query projections at inference to sharpen attention over relevant tokens. However, it only adapts to the given context rather than acquiring new knowledge from user interactions. Similarly, Hu \textit{et al.}~\cite{hu2025test} proposed TLM, which adapts LLMs at test time by minimizing input perplexity via LoRA on high-perplexity samples. However, this primarily reinforces existing predictions rather than incorporating new knowledge. Finally, Tandon \textit{et al.}~\cite{tandon2025end} reframed long-context modeling as continual learning, compressing context into weights via next-token prediction with $\mathcal{O}(1)$ decoding latency. However, this remains contextual compression, as the model does not acquire knowledge beyond the given sequence.

It is worth clarifying that Gao \textit{et al.}~\cite{gao2025can} evaluate reasoning versus memorization through simulated cutoffs rather than unlearning, and update no model weights, unlike both unlearning and TTT approaches discussed above. In summary, existing TTT methods compress context but do not encode new knowledge into model parameters. True test-time learning should enable models to update their knowledge based on user interactions. We demonstrate this through interactive machine unlearning, enabling users to modify model knowledge on-the-fly without requiring training pipelines.